\chapter{Implementaci\'on f\'isica}
\recursoscapitulo{07-implementacion}
El esquema se implementa de la misma manera en ambos gestores: las mismas tablas, columnas, llaves y restricciones, con los mismos nombres. Las diferencias se limitan a la sintaxis y a los tipos de dato propios de cada gestor. Este cap\'itulo describe las decisiones comunes a ambas implementaciones.

\paragraph{Equivalencia de tipos.} El diccionario de datos utiliza tipos genéricos. La tabla~\ref{tab:equivalencia-tipos} muestra su equivalencia en cada gestor.

\begin{table}[H]
	\centering
	\small
	\caption{Equivalencia de tipos de dato entre gestores}
	\label{tab:equivalencia-tipos}
	\begin{tabular}{@{}L{3.3cm} L{3.9cm} L{3.1cm} L{5.5cm}@{}}
		\toprule
		\textbf{Diccionario} & \textbf{PostgreSQL} & \textbf{MySQL} & \textbf{Justificación} \\
		\midrule
		INT y BIGINT autoincremental & \nem{GENERATED BY DEFAULT AS IDENTITY} & \nem{AUTO\us INCREMENT} & El gestor genera la clave, pero admite insertar un valor explícito, lo que permite trasladar los datos a la segunda instancia conservando sus claves. \\
		DECIMAL(p,s) & \nem{NUMERIC(p,s)} & \nem{DECIMAL(p,s)} & Valores exactos, sin errores de redondeo en montos y tasas. \\
		TIMESTAMP & \nem{TIMESTAMP} & \nem{DATETIME} & Fecha y hora local sin conversión. En MySQL, \nem{TIMESTAMP} convierte los valores a UTC y solo admite fechas hasta 2038. \\
		JSON & \nem{JSONB} & \nem{JSON} & Ambos validan que el contenido sea JSON. En PostgreSQL, \nem{JSONB} lo almacena en formato binario, más eficiente para consultarlo. \\
		VARCHAR, CHAR, SMALLINT, DATE, TEXT & Igual & Igual & Tipos estándar con el mismo comportamiento en ambos gestores. \\
		\bottomrule
	\end{tabular}
\end{table}

\paragraph{Nomenclatura de restricciones.} Todas las restricciones se nombran explícitamente con el prefijo de su tipo y la tabla a la que pertenecen: \nem{pk\us} para llaves primarias, \nem{ak\us} para llaves alternas, \nem{fk\us} para llaves foráneas y \nem{ck\us} para reglas de validación. Así, cuando una operación viola una restricción, el mensaje de error identifica la tabla y la regla incumplida, y los nombres son los mismos en ambos gestores.

\paragraph{Validación de las reglas de negocio.} Las reglas que involucran columnas de una sola tabla se implementan como restricciones \nem{CHECK}, que el gestor evalúa en cada inserción y actualización sin código adicional. Las reglas que comparan datos de dos tablas no pueden expresarse con \nem{CHECK}, por lo que se implementan con triggers. La tabla~\ref{tab:validacion-reglas} muestra la distribución.

\begin{table}[H]
	\centering
	\small
	\caption{Mecanismo de validación de cada regla de negocio}
	\label{tab:validacion-reglas}
	\begin{tabular}{@{}L{2.2cm} L{13.6cm}@{}}
		\toprule
		\textbf{Mecanismo} & \textbf{Reglas} \\
		\midrule
		\nem{CHECK} & Capital mayor a cero; tasas no negativas; día de pago entre 1 y 28; vencimiento posterior a la suscripción; vencimiento en el día de pago; fecha final del periodo posterior a la inicial; montos de interés, pago, mora y liquidación no negativos; operación de bitácora válida. \\
		Trigger & RFC del aval distinto del RFC del cliente; capital liquidado igual al invertido; liquidación en la fecha de vencimiento o después. \\
		\bottomrule
	\end{tabular}
\end{table}

El límite superior de las tasas (99.9999) no requiere una restricción, ya que lo impone la precisión del tipo \nem{NUMERIC(6,4)}.

\section{PostgreSQL}

\subsection{Base de datos}

La base de datos la crea el superusuario y la asigna al DBA como dueño, de acuerdo con la sección~\ref{sec:admin-instancia}. Se crea a partir de \nem{template0}, la plantilla sin modificaciones, para garantizar que la codificación UTF-8 se aplique sin heredar objetos de otra plantilla.

\begin{verbatim}
	-- Superusuario
	CREATE DATABASE inversiones
	  OWNER rol_admin_bd
	  ENCODING 'UTF8'
	  TEMPLATE template0;
\end{verbatim}

Dentro de la base de datos, el DBA crea el esquema \nem{inversiones}, donde residen todos los objetos del sistema, y lo establece como ruta de búsqueda predeterminada para que las sentencias no requieran el prefijo del esquema. El esquema \nem{public}, que el gestor crea por defecto, se elimina para que ningún objeto pueda crearse por error fuera del esquema del sistema.

\begin{verbatim}
	-- DBA, conectado a la base de datos inversiones
	SET ROLE rol_admin_bd;
	
	CREATE SCHEMA inversiones AUTHORIZATION rol_admin_bd;
	ALTER DATABASE inversiones SET search_path TO inversiones;
	DROP SCHEMA public;
\end{verbatim}

\begin{verbatim}
	\l inversiones
	\dn
\end{verbatim}

\captura{pg-base-datos.png}{Base de datos y esquema del sistema en PostgreSQL}{pg-base-datos}

\subsection{Tablas}
Las tablas se crean en el orden de sus dependencias: primero las tablas referenciadas y después las que contienen las llaves foráneas. El DBA las crea con el rol \nem{rol\us admin\us bd} activo, de modo que el rol es su dueño.

\begin{verbatim}
	SET ROLE rol_admin_bd;

	CREATE TABLE empresa (
	  empresa_rfc        VARCHAR(13)  NOT NULL,
	  empresa_nombre     VARCHAR(100) NOT NULL,
	  empresa_direccion  VARCHAR(150) NOT NULL,
	  CONSTRAINT pk_empresa PRIMARY KEY (empresa_rfc)
	);

	CREATE TABLE sucursal (
	  sucursal_num        INT GENERATED BY DEFAULT AS IDENTITY,
	  empresa_rfc         VARCHAR(13)  NOT NULL,
	  sucursal_nombre     VARCHAR(60)  NOT NULL,
	  sucursal_ciudad     VARCHAR(100) NOT NULL,
	  sucursal_direccion  VARCHAR(150) NOT NULL,
	  CONSTRAINT pk_sucursal PRIMARY KEY (sucursal_num),
	  CONSTRAINT fk_sucursal_empresa FOREIGN KEY (empresa_rfc)
	    REFERENCES empresa (empresa_rfc)
	    ON DELETE RESTRICT ON UPDATE CASCADE
	);

	CREATE TABLE cliente (
	  cliente_num        INT GENERATED BY DEFAULT AS IDENTITY,
	  cliente_rfc        CHAR(13)     NOT NULL,
	  cliente_nombre     VARCHAR(60)  NOT NULL,
	  cliente_apellido   VARCHAR(80)  NOT NULL,
	  cliente_direccion  VARCHAR(150) NOT NULL,
	  cliente_telefono   VARCHAR(15)  NOT NULL,
	  CONSTRAINT pk_cliente PRIMARY KEY (cliente_num),
	  CONSTRAINT ak_cliente_rfc UNIQUE (cliente_rfc)
	);

	CREATE TABLE aval (
	  aval_num        INT GENERATED BY DEFAULT AS IDENTITY,
	  aval_rfc        CHAR(13)     NOT NULL,
	  aval_nombre     VARCHAR(60)  NOT NULL,
	  aval_apellido   VARCHAR(80)  NOT NULL,
	  aval_direccion  VARCHAR(150) NOT NULL,
	  CONSTRAINT pk_aval PRIMARY KEY (aval_num),
	  CONSTRAINT ak_aval_rfc UNIQUE (aval_rfc)
	);

	CREATE TABLE inversion (
	  inversion_num                INT GENERATED BY DEFAULT AS IDENTITY,
	  pagare_num                   VARCHAR(30)   NOT NULL,
	  sucursal_num                 INT           NOT NULL,
	  cliente_num                  INT           NOT NULL,
	  aval_num                     INT           NULL,
	  inversion_capital            NUMERIC(14,2) NOT NULL,
	  inversion_tasa               NUMERIC(6,4)  NOT NULL,
	  inversion_tasa_moratoria     NUMERIC(6,4)  NOT NULL,
	  inversion_dia_pago           SMALLINT      NOT NULL,
	  pagare_fecha                 DATE          NOT NULL,
	  inversion_moneda             CHAR(3)       NOT NULL DEFAULT 'MXN',
	  pagare_jurisdiccion          VARCHAR(100)  NOT NULL,
	  inversion_fecha_vencimiento  DATE          NOT NULL,
	  CONSTRAINT pk_inversion PRIMARY KEY (inversion_num),
	  CONSTRAINT ak_inversion_pagare UNIQUE (pagare_num),
	  CONSTRAINT fk_inversion_sucursal FOREIGN KEY (sucursal_num)
	    REFERENCES sucursal (sucursal_num)
	    ON DELETE RESTRICT ON UPDATE RESTRICT,
	  CONSTRAINT fk_inversion_cliente FOREIGN KEY (cliente_num)
	    REFERENCES cliente (cliente_num)
	    ON DELETE RESTRICT ON UPDATE RESTRICT,
	  CONSTRAINT fk_inversion_aval FOREIGN KEY (aval_num)
	    REFERENCES aval (aval_num)
	    ON DELETE RESTRICT ON UPDATE RESTRICT,
	  CONSTRAINT ck_inversion_capital CHECK (inversion_capital > 0),
	  CONSTRAINT ck_inversion_tasa CHECK (inversion_tasa >= 0),
	  CONSTRAINT ck_inversion_tasa_moratoria
	    CHECK (inversion_tasa_moratoria >= 0),
	  CONSTRAINT ck_inversion_dia_pago
	    CHECK (inversion_dia_pago BETWEEN 1 AND 28),
	  CONSTRAINT ck_inversion_vencimiento
	    CHECK (inversion_fecha_vencimiento > pagare_fecha),
	  CONSTRAINT ck_inversion_vencimiento_dia
	    CHECK (EXTRACT(DAY FROM inversion_fecha_vencimiento)
	           = inversion_dia_pago)
	);

	CREATE TABLE interes (
	  inversion_num          INT           NOT NULL,
	  interes_fecha_inicial  DATE          NOT NULL,
	  interes_fecha_final    DATE          NOT NULL,
	  interes_monto          NUMERIC(14,2) NOT NULL,
	  CONSTRAINT pk_interes PRIMARY KEY (inversion_num, interes_fecha_inicial),
	  CONSTRAINT fk_interes_inversion FOREIGN KEY (inversion_num)
	    REFERENCES inversion (inversion_num)
	    ON DELETE RESTRICT ON UPDATE RESTRICT,
	  CONSTRAINT ck_interes_fechas
	    CHECK (interes_fecha_final > interes_fecha_inicial),
	  CONSTRAINT ck_interes_monto CHECK (interes_monto >= 0)
	);

	CREATE TABLE pago (
	  inversion_num          INT           NOT NULL,
	  interes_fecha_inicial  DATE          NOT NULL,
	  pago_fecha             DATE          NOT NULL,
	  pago_interes           NUMERIC(14,2) NOT NULL,
	  pago_mora              NUMERIC(14,2) NOT NULL DEFAULT 0,
	  CONSTRAINT pk_pago
	    PRIMARY KEY (inversion_num, interes_fecha_inicial, pago_fecha),
	  CONSTRAINT fk_pago_interes FOREIGN KEY (inversion_num, interes_fecha_inicial)
	    REFERENCES interes (inversion_num, interes_fecha_inicial)
	    ON DELETE RESTRICT ON UPDATE RESTRICT,
	  CONSTRAINT ck_pago_montos CHECK (pago_interes >= 0 AND pago_mora >= 0)
	);

	CREATE TABLE liquidacion (
	  inversion_num        INT           NOT NULL,
	  liquidacion_fecha    DATE          NOT NULL,
	  liquidacion_capital  NUMERIC(14,2) NOT NULL,
	  liquidacion_mora     NUMERIC(14,2) NOT NULL DEFAULT 0,
	  CONSTRAINT pk_liquidacion PRIMARY KEY (inversion_num),
	  CONSTRAINT fk_liquidacion_inversion FOREIGN KEY (inversion_num)
	    REFERENCES inversion (inversion_num)
	    ON DELETE RESTRICT ON UPDATE RESTRICT,
	  CONSTRAINT ck_liquidacion_capital CHECK (liquidacion_capital > 0),
	  CONSTRAINT ck_liquidacion_mora CHECK (liquidacion_mora >= 0)
	);

	CREATE TABLE bitacora (
	  bitacora_num            BIGINT GENERATED BY DEFAULT AS IDENTITY,
	  bitacora_fecha          TIMESTAMP    NOT NULL DEFAULT LOCALTIMESTAMP,
	  bitacora_usuario        VARCHAR(100) NOT NULL,
	  bitacora_tabla          VARCHAR(63)  NOT NULL,
	  bitacora_operacion      VARCHAR(6)   NOT NULL,
	  bitacora_ip             VARCHAR(45)  NULL,
	  bitacora_sesion         BIGINT       NULL,
	  bitacora_query          TEXT         NULL,
	  bitacora_dato_anterior  JSONB        NULL,
	  bitacora_dato_nuevo     JSONB        NULL,
	  CONSTRAINT pk_bitacora PRIMARY KEY (bitacora_num, bitacora_fecha),
	  CONSTRAINT ck_bitacora_operacion
	    CHECK (bitacora_operacion IN ('INSERT', 'UPDATE', 'DELETE'))
	);
\end{verbatim}

La fecha de la bitácora se asigna con \nem{LOCALTIMESTAMP}, que devuelve la hora local sin zona horaria, en lugar de \nem{CURRENT\us TIMESTAMP}, que la incluye y obligaría a convertirla para almacenarla en una columna \nem{TIMESTAMP}.

Las tablas y sus restricciones se consultan en el catálogo del gestor (figura~\ref{fig:pg-tablas}):

\begin{verbatim}
	SELECT conrelid::regclass AS tabla, conname AS restriccion,
	       CASE contype WHEN 'p' THEN 'PK' WHEN 'u' THEN 'AK'
	                    WHEN 'f' THEN 'FK' WHEN 'c' THEN 'CHECK' END AS tipo
	FROM pg_constraint
	WHERE connamespace = 'inversiones'::regnamespace
	ORDER BY 1, 3, 2;
\end{verbatim}

\captura{pg-tablas.png}{Tablas y restricciones del esquema en PostgreSQL}{pg-tablas}

\subsection{Objetos programados}

\paragraph{Procedimiento \nem{generar\us tabla\us interes}.} Genera los periodos de interés de una inversión a partir de las condiciones de su pagaré. Se implementa como procedimiento almacenado, y no como un cálculo que realice cada usuario, para que todas las tablas de interés se generen con la misma regla:

\begin{itemize}[leftmargin=*]
	\item El monto de cada periodo es fijo e igual a capital $\times$ tasa / 100, ya que el pagaré establece un importe mensual fijo.
	\item El primer periodo inicia en la fecha de suscripción y termina en el día de pago del mes siguiente; los demás van de un día de pago al siguiente.
	\item El último periodo termina en la fecha de vencimiento, que por regla coincide con el día de pago.
\end{itemize}

El procedimiento rechaza una inversión que ya tenga tabla de interés, para no duplicar periodos, y una cuyo vencimiento sea anterior al primer día de pago, que no tendría ningún periodo completo.

Se define con \nem{SECURITY DEFINER}, es decir, se ejecuta con los privilegios de su dueño. De esta forma, el ejecutivo puede generar los periodos sin tener privilegio de inserción sobre \nem{interes}. Por seguridad, el procedimiento fija su propia ruta de búsqueda, de modo que no puede ser engañado con objetos del mismo nombre creados en otro esquema.

\begin{verbatim}
	CREATE PROCEDURE generar_tabla_interes(p_inversion_num INT)
	LANGUAGE plpgsql
	SECURITY DEFINER
	SET search_path = inversiones
	AS $$
	DECLARE
	  v        inversion%ROWTYPE;
	  v_monto  NUMERIC(14,2);
	  v_inicio DATE;
	  v_fin    DATE;
	BEGIN
	  SELECT * INTO v FROM inversion WHERE inversion_num = p_inversion_num;
	  IF NOT FOUND THEN
	    RAISE EXCEPTION 'La inversión % no existe', p_inversion_num;
	  END IF;

	  IF EXISTS (SELECT 1 FROM interes WHERE inversion_num = p_inversion_num) THEN
	    RAISE EXCEPTION 'La inversión % ya tiene tabla de interés',
	      p_inversion_num;
	  END IF;

	  v_monto  := round(v.inversion_capital * v.inversion_tasa / 100, 2);
	  v_inicio := v.pagare_fecha;
	  v_fin    := (date_trunc('month', v.pagare_fecha) + INTERVAL '1 month'
	               + (v.inversion_dia_pago - 1) * INTERVAL '1 day')::DATE;

	  IF v_fin > v.inversion_fecha_vencimiento THEN
	    RAISE EXCEPTION 'El vencimiento es anterior al primer día de pago';
	  END IF;

	  WHILE v_fin <= v.inversion_fecha_vencimiento LOOP
	    INSERT INTO interes (inversion_num, interes_fecha_inicial,
	                         interes_fecha_final, interes_monto)
	    VALUES (p_inversion_num, v_inicio, v_fin, v_monto);
	    v_inicio := v_fin;
	    v_fin    := (v_fin + INTERVAL '1 month')::DATE;
	  END LOOP;
	END;
	$$;
\end{verbatim}

Como el día de pago es como máximo 28, sumar un mes a la fecha nunca produce un día inexistente, como el 30 de febrero.

\paragraph{Triggers de validación.} Las reglas que comparan dos tablas se validan antes de cada inserción o actualización. Si la regla no se cumple, el trigger cancela la operación con un mensaje que explica el motivo. A diferencia del procedimiento, estas funciones se ejecutan con los privilegios del usuario que realiza la operación, ya que solo consultan tablas que el ejecutivo y la tesorería pueden leer.

\begin{verbatim}
	CREATE FUNCTION fn_validar_aval() RETURNS trigger
	LANGUAGE plpgsql
	SET search_path = inversiones
	AS $$
	BEGIN
	  IF NEW.aval_num IS NOT NULL AND
	     (SELECT aval_rfc FROM aval WHERE aval_num = NEW.aval_num) =
	     (SELECT cliente_rfc FROM cliente WHERE cliente_num = NEW.cliente_num)
	  THEN
	    RAISE EXCEPTION 'El aval no puede ser el cliente de la inversión';
	  END IF;
	  RETURN NEW;
	END;
	$$;

	CREATE TRIGGER trg_inversion_aval
	  BEFORE INSERT OR UPDATE OF aval_num, cliente_num ON inversion
	  FOR EACH ROW EXECUTE FUNCTION fn_validar_aval();

	CREATE FUNCTION fn_validar_liquidacion() RETURNS trigger
	LANGUAGE plpgsql
	SET search_path = inversiones
	AS $$
	DECLARE
	  v_capital      NUMERIC(14,2);
	  v_vencimiento  DATE;
	BEGIN
	  SELECT inversion_capital, inversion_fecha_vencimiento
	    INTO v_capital, v_vencimiento
	  FROM inversion
	  WHERE inversion_num = NEW.inversion_num;

	  IF NEW.liquidacion_capital <> v_capital THEN
	    RAISE EXCEPTION 'El capital liquidado debe ser igual al invertido (%)',
	      v_capital;
	  END IF;

	  IF NEW.liquidacion_fecha < v_vencimiento THEN
	    RAISE EXCEPTION 'La liquidación no puede ser anterior al vencimiento (%)',
	      v_vencimiento;
	  END IF;
	  RETURN NEW;
	END;
	$$;

	CREATE TRIGGER trg_liquidacion_validar
	  BEFORE INSERT OR UPDATE ON liquidacion
	  FOR EACH ROW EXECUTE FUNCTION fn_validar_liquidacion();
\end{verbatim}

\begin{verbatim}
	SELECT p.proname AS objeto,
	       CASE p.prokind WHEN 'p' THEN 'Procedimiento' ELSE 'Función' END AS tipo,
	       pg_get_userbyid(p.proowner) AS dueno,
	       p.prosecdef AS security_definer
	FROM pg_proc p
	WHERE p.pronamespace = 'inversiones'::regnamespace;

	SELECT tgname AS trigger, tgrelid::regclass AS tabla
	FROM pg_trigger
	WHERE NOT tgisinternal;
\end{verbatim}

\captura{pg-objetos.png}{Procedimiento, funciones y triggers en PostgreSQL}{pg-objetos}

\subsection{Permisos}

Los privilegios de la matriz de la tabla~\ref{tab:matriz} se otorgan a los roles, nunca a los usuarios, de modo que un cambio de puesto solo requiere cambiar el rol de la persona. Antes de otorgarlos, se retiran los privilegios que PostgreSQL concede por defecto a todos los usuarios (\nem{PUBLIC}): la conexión a la base de datos y la ejecución de funciones y procedimientos. Sin este paso, cualquier usuario de la instancia podría conectarse a la base de datos o ejecutar el procedimiento, aunque no tuviera un rol del sistema.

\begin{verbatim}
	SET ROLE rol_admin_bd;

	-- Privilegios por defecto de PUBLIC
	REVOKE CONNECT, TEMPORARY ON DATABASE inversiones FROM PUBLIC;
	REVOKE EXECUTE ON PROCEDURE generar_tabla_interes(INT) FROM PUBLIC;
	REVOKE EXECUTE ON FUNCTION fn_validar_aval(), fn_validar_liquidacion()
	  FROM PUBLIC;

	-- Acceso a la base de datos y al esquema
	GRANT CONNECT ON DATABASE inversiones
	  TO rol_gerente, rol_ejecutivo, rol_tesoreria, rol_auditor, rol_respaldo;
	GRANT USAGE ON SCHEMA inversiones
	  TO rol_gerente, rol_ejecutivo, rol_tesoreria, rol_auditor, rol_respaldo;

	-- Gerente
	GRANT SELECT, UPDATE         ON empresa, cliente           TO rol_gerente;
	GRANT SELECT, INSERT, UPDATE ON sucursal, aval             TO rol_gerente;
	GRANT SELECT, UPDATE, DELETE ON inversion                  TO rol_gerente;
	GRANT SELECT                 ON interes, pago, liquidacion TO rol_gerente;
	GRANT USAGE ON SEQUENCE sucursal_sucursal_num_seq, aval_aval_num_seq
	  TO rol_gerente;

	-- Ejecutivo
	GRANT SELECT ON empresa, sucursal, aval, interes, pago, liquidacion
	  TO rol_ejecutivo;
	GRANT SELECT, INSERT ON cliente, inversion TO rol_ejecutivo;
	GRANT UPDATE (cliente_telefono, cliente_direccion) ON cliente
	  TO rol_ejecutivo;
	GRANT USAGE ON SEQUENCE cliente_cliente_num_seq, inversion_inversion_num_seq
	  TO rol_ejecutivo;
	GRANT EXECUTE ON PROCEDURE generar_tabla_interes(INT) TO rol_ejecutivo;

	-- Tesorería
	GRANT SELECT ON empresa, sucursal, cliente, aval, inversion, interes
	  TO rol_tesoreria;
	GRANT SELECT, INSERT ON pago, liquidacion TO rol_tesoreria;

	-- Auditor y respaldo
	GRANT SELECT ON ALL TABLES IN SCHEMA inversiones TO rol_auditor, rol_respaldo;
	GRANT SELECT ON ALL SEQUENCES IN SCHEMA inversiones TO rol_respaldo;

	-- Tablas creadas posteriormente
	ALTER DEFAULT PRIVILEGES FOR ROLE rol_admin_bd IN SCHEMA inversiones
	  GRANT SELECT ON TABLES TO rol_auditor, rol_respaldo;
\end{verbatim}

El privilegio \nem{USAGE} sobre las secuencias de las columnas de identidad permite a los roles que insertan registros generar sus claves. El último bloque establece privilegios por defecto: toda tabla que el DBA cree después, como las particiones de la segunda instancia, será legible automáticamente por el auditor y por el usuario de respaldo, sin necesidad de otorgar el privilegio en cada ocasión.

Los privilegios otorgados se consultan con el comando \nem{\textbackslash dp} (figura~\ref{fig:pg-permisos}):

\begin{verbatim}
	\dp inversiones.*
\end{verbatim}

\captura{pg-permisos.png}{Privilegios sobre las tablas en PostgreSQL}{pg-permisos}

\section{MySQL}

\subsection{Base de datos}

En MySQL, la base de datos y el esquema son el mismo objeto, por lo que no se crea un esquema adicional. El DBA crea la base de datos, ya que recibió todos los privilegios sobre \nem{inversiones} desde antes de que existiera (sección~\ref{sec:roles-mysql}).

La base de datos utiliza la codificación \nem{utf8mb4} y la intercalación \nem{utf8mb4\us 0900\us ai\us ci}, que compara los textos sin distinguir mayúsculas, minúsculas ni acentos. Así, una búsqueda por nombre encuentra a ``Jose'' aunque esté registrado como ``José''. PostgreSQL, en cambio, distingue mayúsculas y acentos; las consultas por nombre de ambos gestores se escriben de modo que devuelvan los mismos resultados (sección~\ref{sec:consultas-objetivos}).

\begin{verbatim}
	-- DBA
	CREATE DATABASE inversiones
	  CHARACTER SET utf8mb4
	  COLLATE utf8mb4_0900_ai_ci;
\end{verbatim}

\begin{verbatim}
	SELECT SCHEMA_NAME, DEFAULT_CHARACTER_SET_NAME, DEFAULT_COLLATION_NAME
	FROM information_schema.SCHEMATA
	WHERE SCHEMA_NAME = 'inversiones';
\end{verbatim}

\captura{mysql-base-datos.png}{Base de datos del sistema en MySQL}{mysql-base-datos}

\subsection{Tablas}

Las tablas tienen las mismas columnas, llaves y restricciones que en PostgreSQL, con los tipos equivalentes de la tabla~\ref{tab:equivalencia-tipos}. Todas utilizan el motor InnoDB, el único de MySQL que admite llaves foráneas y transacciones; aunque es el motor por defecto, se indica de forma explícita para que la definición no dependa de la configuración del servidor.

\begin{verbatim}
	USE inversiones;

	CREATE TABLE empresa (
	  empresa_rfc        VARCHAR(13)  NOT NULL,
	  empresa_nombre     VARCHAR(100) NOT NULL,
	  empresa_direccion  VARCHAR(150) NOT NULL,
	  CONSTRAINT pk_empresa PRIMARY KEY (empresa_rfc)
	) ENGINE = InnoDB;

	CREATE TABLE sucursal (
	  sucursal_num        INT          NOT NULL AUTO_INCREMENT,
	  empresa_rfc         VARCHAR(13)  NOT NULL,
	  sucursal_nombre     VARCHAR(60)  NOT NULL,
	  sucursal_ciudad     VARCHAR(100) NOT NULL,
	  sucursal_direccion  VARCHAR(150) NOT NULL,
	  CONSTRAINT pk_sucursal PRIMARY KEY (sucursal_num),
	  CONSTRAINT fk_sucursal_empresa FOREIGN KEY (empresa_rfc)
	    REFERENCES empresa (empresa_rfc)
	    ON DELETE RESTRICT ON UPDATE CASCADE
	) ENGINE = InnoDB;

	CREATE TABLE cliente (
	  cliente_num        INT          NOT NULL AUTO_INCREMENT,
	  cliente_rfc        CHAR(13)     NOT NULL,
	  cliente_nombre     VARCHAR(60)  NOT NULL,
	  cliente_apellido   VARCHAR(80)  NOT NULL,
	  cliente_direccion  VARCHAR(150) NOT NULL,
	  cliente_telefono   VARCHAR(15)  NOT NULL,
	  CONSTRAINT pk_cliente PRIMARY KEY (cliente_num),
	  CONSTRAINT ak_cliente_rfc UNIQUE (cliente_rfc)
	) ENGINE = InnoDB;

	CREATE TABLE aval (
	  aval_num        INT          NOT NULL AUTO_INCREMENT,
	  aval_rfc        CHAR(13)     NOT NULL,
	  aval_nombre     VARCHAR(60)  NOT NULL,
	  aval_apellido   VARCHAR(80)  NOT NULL,
	  aval_direccion  VARCHAR(150) NOT NULL,
	  CONSTRAINT pk_aval PRIMARY KEY (aval_num),
	  CONSTRAINT ak_aval_rfc UNIQUE (aval_rfc)
	) ENGINE = InnoDB;

	CREATE TABLE inversion (
	  inversion_num                INT           NOT NULL AUTO_INCREMENT,
	  pagare_num                   VARCHAR(30)   NOT NULL,
	  sucursal_num                 INT           NOT NULL,
	  cliente_num                  INT           NOT NULL,
	  aval_num                     INT           NULL,
	  inversion_capital            DECIMAL(14,2) NOT NULL,
	  inversion_tasa               DECIMAL(6,4)  NOT NULL,
	  inversion_tasa_moratoria     DECIMAL(6,4)  NOT NULL,
	  inversion_dia_pago           SMALLINT      NOT NULL,
	  pagare_fecha                 DATE          NOT NULL,
	  inversion_moneda             CHAR(3)       NOT NULL DEFAULT 'MXN',
	  pagare_jurisdiccion          VARCHAR(100)  NOT NULL,
	  inversion_fecha_vencimiento  DATE          NOT NULL,
	  CONSTRAINT pk_inversion PRIMARY KEY (inversion_num),
	  CONSTRAINT ak_inversion_pagare UNIQUE (pagare_num),
	  CONSTRAINT fk_inversion_sucursal FOREIGN KEY (sucursal_num)
	    REFERENCES sucursal (sucursal_num)
	    ON DELETE RESTRICT ON UPDATE RESTRICT,
	  CONSTRAINT fk_inversion_cliente FOREIGN KEY (cliente_num)
	    REFERENCES cliente (cliente_num)
	    ON DELETE RESTRICT ON UPDATE RESTRICT,
	  CONSTRAINT fk_inversion_aval FOREIGN KEY (aval_num)
	    REFERENCES aval (aval_num)
	    ON DELETE RESTRICT ON UPDATE RESTRICT,
	  CONSTRAINT ck_inversion_capital CHECK (inversion_capital > 0),
	  CONSTRAINT ck_inversion_tasa CHECK (inversion_tasa >= 0),
	  CONSTRAINT ck_inversion_tasa_moratoria
	    CHECK (inversion_tasa_moratoria >= 0),
	  CONSTRAINT ck_inversion_dia_pago
	    CHECK (inversion_dia_pago BETWEEN 1 AND 28),
	  CONSTRAINT ck_inversion_vencimiento
	    CHECK (inversion_fecha_vencimiento > pagare_fecha),
	  CONSTRAINT ck_inversion_vencimiento_dia
	    CHECK (DAY(inversion_fecha_vencimiento) = inversion_dia_pago)
	) ENGINE = InnoDB;

	CREATE TABLE interes (
	  inversion_num          INT           NOT NULL,
	  interes_fecha_inicial  DATE          NOT NULL,
	  interes_fecha_final    DATE          NOT NULL,
	  interes_monto          DECIMAL(14,2) NOT NULL,
	  CONSTRAINT pk_interes PRIMARY KEY (inversion_num, interes_fecha_inicial),
	  CONSTRAINT fk_interes_inversion FOREIGN KEY (inversion_num)
	    REFERENCES inversion (inversion_num)
	    ON DELETE RESTRICT ON UPDATE RESTRICT,
	  CONSTRAINT ck_interes_fechas
	    CHECK (interes_fecha_final > interes_fecha_inicial),
	  CONSTRAINT ck_interes_monto CHECK (interes_monto >= 0)
	) ENGINE = InnoDB;

	CREATE TABLE pago (
	  inversion_num          INT           NOT NULL,
	  interes_fecha_inicial  DATE          NOT NULL,
	  pago_fecha             DATE          NOT NULL,
	  pago_interes           DECIMAL(14,2) NOT NULL,
	  pago_mora              DECIMAL(14,2) NOT NULL DEFAULT 0,
	  CONSTRAINT pk_pago
	    PRIMARY KEY (inversion_num, interes_fecha_inicial, pago_fecha),
	  CONSTRAINT fk_pago_interes FOREIGN KEY (inversion_num, interes_fecha_inicial)
	    REFERENCES interes (inversion_num, interes_fecha_inicial)
	    ON DELETE RESTRICT ON UPDATE RESTRICT,
	  CONSTRAINT ck_pago_montos CHECK (pago_interes >= 0 AND pago_mora >= 0)
	) ENGINE = InnoDB;

	CREATE TABLE liquidacion (
	  inversion_num        INT           NOT NULL,
	  liquidacion_fecha    DATE          NOT NULL,
	  liquidacion_capital  DECIMAL(14,2) NOT NULL,
	  liquidacion_mora     DECIMAL(14,2) NOT NULL DEFAULT 0,
	  CONSTRAINT pk_liquidacion PRIMARY KEY (inversion_num),
	  CONSTRAINT fk_liquidacion_inversion FOREIGN KEY (inversion_num)
	    REFERENCES inversion (inversion_num)
	    ON DELETE RESTRICT ON UPDATE RESTRICT,
	  CONSTRAINT ck_liquidacion_capital CHECK (liquidacion_capital > 0),
	  CONSTRAINT ck_liquidacion_mora CHECK (liquidacion_mora >= 0)
	) ENGINE = InnoDB;

	CREATE TABLE bitacora (
	  bitacora_num            BIGINT       NOT NULL AUTO_INCREMENT,
	  bitacora_fecha          DATETIME     NOT NULL DEFAULT CURRENT_TIMESTAMP,
	  bitacora_usuario        VARCHAR(100) NOT NULL,
	  bitacora_tabla          VARCHAR(63)  NOT NULL,
	  bitacora_operacion      VARCHAR(6)   NOT NULL,
	  bitacora_ip             VARCHAR(45)  NULL,
	  bitacora_sesion         BIGINT       NULL,
	  bitacora_query          TEXT         NULL,
	  bitacora_dato_anterior  JSON         NULL,
	  bitacora_dato_nuevo     JSON         NULL,
	  CONSTRAINT pk_bitacora PRIMARY KEY (bitacora_num, bitacora_fecha),
	  CONSTRAINT ck_bitacora_operacion
	    CHECK (bitacora_operacion IN ('INSERT', 'UPDATE', 'DELETE'))
	) ENGINE = InnoDB;
\end{verbatim}

En MySQL, \nem{CURRENT\us TIMESTAMP} asignado a una columna \nem{DATETIME} devuelve la hora local según la zona horaria del servidor (\nem{default\us time\us zone}), por lo que cumple la misma función que \nem{LOCALTIMESTAMP} en PostgreSQL.

Las restricciones de cada tabla se consultan en el diccionario de datos del gestor (figura~\ref{fig:mysql-tablas}):

\begin{verbatim}
	SELECT TABLE_NAME, CONSTRAINT_NAME, CONSTRAINT_TYPE
	FROM information_schema.TABLE_CONSTRAINTS
	WHERE CONSTRAINT_SCHEMA = 'inversiones'
	ORDER BY TABLE_NAME, CONSTRAINT_TYPE, CONSTRAINT_NAME;
\end{verbatim}

\captura{mysql-tablas.png}{Tablas y restricciones del esquema en MySQL}{mysql-tablas}

\subsection{Objetos programados}

\paragraph{Procedimiento \nem{generar\us tabla\us interes}.} Aplica la misma regla que en PostgreSQL. Se define con \nem{SQL SECURITY DEFINER}, equivalente a \nem{SECURITY DEFINER}, y su \nem{DEFINER} es la cuenta del DBA que lo crea, de acuerdo con la sección~\ref{sec:admin-instancia}.

\begin{verbatim}
	DELIMITER //

	CREATE DEFINER = CURRENT_USER PROCEDURE generar_tabla_interes(
	  IN p_inversion_num INT)
	SQL SECURITY DEFINER
	MODIFIES SQL DATA
	BEGIN
	  DECLARE v_capital     DECIMAL(14,2);
	  DECLARE v_tasa        DECIMAL(6,4);
	  DECLARE v_dia         SMALLINT;
	  DECLARE v_suscripcion DATE;
	  DECLARE v_vencimiento DATE;
	  DECLARE v_monto       DECIMAL(14,2);
	  DECLARE v_inicio      DATE;
	  DECLARE v_fin         DATE;

	  SELECT inversion_capital, inversion_tasa, inversion_dia_pago,
	         pagare_fecha, inversion_fecha_vencimiento
	    INTO v_capital, v_tasa, v_dia, v_suscripcion, v_vencimiento
	  FROM inversion
	  WHERE inversion_num = p_inversion_num;

	  IF v_capital IS NULL THEN
	    SIGNAL SQLSTATE '45000'
	      SET MESSAGE_TEXT = 'La inversión no existe';
	  END IF;

	  IF EXISTS (SELECT 1 FROM interes
	             WHERE inversion_num = p_inversion_num) THEN
	    SIGNAL SQLSTATE '45000'
	      SET MESSAGE_TEXT = 'La inversión ya tiene tabla de interés';
	  END IF;

	  SET v_monto  = ROUND(v_capital * v_tasa / 100, 2);
	  SET v_inicio = v_suscripcion;
	  -- Primer día del mes siguiente, más los días hasta el día de pago
	  SET v_fin = DATE_ADD(
	                DATE_ADD(DATE_SUB(v_suscripcion,
	                                  INTERVAL DAY(v_suscripcion) - 1 DAY),
	                         INTERVAL 1 MONTH),
	                INTERVAL v_dia - 1 DAY);

	  IF v_fin > v_vencimiento THEN
	    SIGNAL SQLSTATE '45000'
	      SET MESSAGE_TEXT = 'El vencimiento es anterior al primer día de pago';
	  END IF;

	  WHILE v_fin <= v_vencimiento DO
	    INSERT INTO interes (inversion_num, interes_fecha_inicial,
	                         interes_fecha_final, interes_monto)
	    VALUES (p_inversion_num, v_inicio, v_fin, v_monto);
	    SET v_inicio = v_fin;
	    SET v_fin    = DATE_ADD(v_fin, INTERVAL 1 MONTH);
	  END WHILE;
	END //

	DELIMITER ;
\end{verbatim}

\paragraph{Validaciones entre tablas.} MySQL no permite que un trigger responda a más de un tipo de operación, por lo que cada regla requiere un trigger para la inserción y otro para la actualización. Para no duplicar la lógica, cada regla se escribe una sola vez en un procedimiento, y ambos triggers lo invocan. A diferencia de PostgreSQL, los triggers de MySQL siempre se ejecutan con los privilegios de su \nem{DEFINER}; como estos procedimientos solo consultan datos, esto no otorga a los usuarios ninguna capacidad adicional.

\begin{verbatim}
	DELIMITER //

	CREATE DEFINER = CURRENT_USER PROCEDURE validar_aval(
	  IN p_cliente_num INT, IN p_aval_num INT)
	READS SQL DATA
	BEGIN
	  IF p_aval_num IS NOT NULL AND
	     (SELECT aval_rfc FROM aval WHERE aval_num = p_aval_num) =
	     (SELECT cliente_rfc FROM cliente WHERE cliente_num = p_cliente_num)
	  THEN
	    SIGNAL SQLSTATE '45000'
	      SET MESSAGE_TEXT = 'El aval no puede ser el cliente de la inversión';
	  END IF;
	END //

	CREATE DEFINER = CURRENT_USER PROCEDURE validar_liquidacion(
	  IN p_inversion_num INT, IN p_fecha DATE, IN p_capital DECIMAL(14,2))
	READS SQL DATA
	BEGIN
	  DECLARE v_capital     DECIMAL(14,2);
	  DECLARE v_vencimiento DATE;

	  SELECT inversion_capital, inversion_fecha_vencimiento
	    INTO v_capital, v_vencimiento
	  FROM inversion
	  WHERE inversion_num = p_inversion_num;

	  IF p_capital <> v_capital THEN
	    SIGNAL SQLSTATE '45000'
	      SET MESSAGE_TEXT = 'El capital liquidado debe ser igual al invertido';
	  END IF;

	  IF p_fecha < v_vencimiento THEN
	    SIGNAL SQLSTATE '45000'
	      SET MESSAGE_TEXT = 'La liquidación no puede ser anterior al vencimiento';
	  END IF;
	END //

	CREATE DEFINER = CURRENT_USER TRIGGER trg_inversion_aval_ins
	  BEFORE INSERT ON inversion FOR EACH ROW
	BEGIN
	  CALL validar_aval(NEW.cliente_num, NEW.aval_num);
	END //

	CREATE DEFINER = CURRENT_USER TRIGGER trg_inversion_aval_upd
	  BEFORE UPDATE ON inversion FOR EACH ROW
	BEGIN
	  CALL validar_aval(NEW.cliente_num, NEW.aval_num);
	END //

	CREATE DEFINER = CURRENT_USER TRIGGER trg_liquidacion_validar_ins
	  BEFORE INSERT ON liquidacion FOR EACH ROW
	BEGIN
	  CALL validar_liquidacion(NEW.inversion_num, NEW.liquidacion_fecha,
	                           NEW.liquidacion_capital);
	END //

	CREATE DEFINER = CURRENT_USER TRIGGER trg_liquidacion_validar_upd
	  BEFORE UPDATE ON liquidacion FOR EACH ROW
	BEGIN
	  CALL validar_liquidacion(NEW.inversion_num, NEW.liquidacion_fecha,
	                           NEW.liquidacion_capital);
	END //

	DELIMITER ;
\end{verbatim}

\begin{verbatim}
	SELECT ROUTINE_NAME, ROUTINE_TYPE, DEFINER, SECURITY_TYPE
	FROM information_schema.ROUTINES
	WHERE ROUTINE_SCHEMA = 'inversiones';
	
	SELECT TRIGGER_NAME, EVENT_MANIPULATION, EVENT_OBJECT_TABLE, DEFINER
	FROM information_schema.TRIGGERS
	WHERE TRIGGER_SCHEMA = 'inversiones';
\end{verbatim}

\captura{mysql-objetos.png}{Procedimientos y triggers en MySQL}{mysql-objetos}

\subsection{Permisos}

Los privilegios se otorgan a los roles, con la misma matriz que en PostgreSQL. Hay dos diferencias que simplifican la implementación:

\begin{itemize}[leftmargin=*]
	\item MySQL no concede privilegios por defecto a todos los usuarios, por lo que no es necesario retirar privilegios antes de otorgar los de cada rol. Un usuario sin privilegios sobre \nem{inversiones} puede iniciar sesión en el servidor, pero no puede ver ni usar la base de datos.
	\item No existen secuencias: \nem{AUTO\us INCREMENT} genera las claves sin requerir un privilegio adicional.
\end{itemize}

Los privilegios del auditor se otorgan sobre la base de datos completa (\nem{inversiones.*}) y no tabla por tabla, con lo que se aplican también a las tablas que se creen después. Este es el equivalente de los privilegios por defecto de PostgreSQL. El usuario de respaldo ya cuenta con sus privilegios sobre la base de datos (sección~\ref{sec:roles-mysql}).

\begin{verbatim}
	-- Gerente
	GRANT SELECT, UPDATE         ON inversiones.empresa     TO rol_gerente;
	GRANT SELECT, INSERT, UPDATE ON inversiones.sucursal    TO rol_gerente;
	GRANT SELECT, UPDATE         ON inversiones.cliente     TO rol_gerente;
	GRANT SELECT, INSERT, UPDATE ON inversiones.aval        TO rol_gerente;
	GRANT SELECT, UPDATE, DELETE ON inversiones.inversion   TO rol_gerente;
	GRANT SELECT                 ON inversiones.interes     TO rol_gerente;
	GRANT SELECT                 ON inversiones.pago        TO rol_gerente;
	GRANT SELECT                 ON inversiones.liquidacion TO rol_gerente;

	-- Ejecutivo
	GRANT SELECT         ON inversiones.empresa     TO rol_ejecutivo;
	GRANT SELECT         ON inversiones.sucursal    TO rol_ejecutivo;
	GRANT SELECT, INSERT ON inversiones.cliente     TO rol_ejecutivo;
	GRANT UPDATE (cliente_telefono, cliente_direccion)
	  ON inversiones.cliente TO rol_ejecutivo;
	GRANT SELECT         ON inversiones.aval        TO rol_ejecutivo;
	GRANT SELECT, INSERT ON inversiones.inversion   TO rol_ejecutivo;
	GRANT SELECT         ON inversiones.interes     TO rol_ejecutivo;
	GRANT SELECT         ON inversiones.pago        TO rol_ejecutivo;
	GRANT SELECT         ON inversiones.liquidacion TO rol_ejecutivo;
	GRANT EXECUTE ON PROCEDURE inversiones.generar_tabla_interes
	  TO rol_ejecutivo;

	-- Tesorería
	GRANT SELECT         ON inversiones.empresa     TO rol_tesoreria;
	GRANT SELECT         ON inversiones.sucursal    TO rol_tesoreria;
	GRANT SELECT         ON inversiones.cliente     TO rol_tesoreria;
	GRANT SELECT         ON inversiones.aval        TO rol_tesoreria;
	GRANT SELECT         ON inversiones.inversion   TO rol_tesoreria;
	GRANT SELECT         ON inversiones.interes     TO rol_tesoreria;
	GRANT SELECT, INSERT ON inversiones.pago        TO rol_tesoreria;
	GRANT SELECT, INSERT ON inversiones.liquidacion TO rol_tesoreria;

	-- Auditor
	GRANT SELECT ON inversiones.* TO rol_auditor;
\end{verbatim}

Los privilegios de cada rol se consultan con \nem{SHOW GRANTS} (figura~\ref{fig:mysql-permisos}):

\begin{verbatim}
	SHOW GRANTS FOR rol_gerente;
	SHOW GRANTS FOR rol_ejecutivo;
	SHOW GRANTS FOR rol_tesoreria;
	SHOW GRANTS FOR rol_auditor;
	SHOW GRANTS FOR rol_respaldo;
\end{verbatim}

\captura{mysql-permisos.png}{Privilegios de los roles en MySQL}{mysql-permisos}

\section{Datos de prueba}

Los datos de prueba son ficticios y representan la operación de la empresa entre enero de 2024 y la fecha de carga, el 28 de octubre de 2026. El volumen supera el mínimo de quince registros por tabla por dos razones. Primero, cada estado de una inversión y cada forma de pago debe tener varios registros, para que las consultas de los objetivos devuelvan resultados significativos. Segundo, las tablas que se particionan en la segunda instancia deben tener registros en cada rango. Aun así, el volumen se mantiene pequeño para que la carga ocupe una fracción mínima de la memoria asignada a cada instancia.

La tabla~\ref{tab:escenarios-prueba} describe los escenarios que cubren los datos y la tabla~\ref{tab:volumen-prueba} el número de registros por tabla.

\begin{table}[H]
	\centering
	\small
	\caption{Escenarios cubiertos por los datos de prueba}
	\label{tab:escenarios-prueba}
	\begin{tabular}{@{}L{4.2cm} L{3.4cm} L{8.4cm}@{}}
		\toprule
		\textbf{Escenario} & \textbf{Registros} & \textbf{Propósito} \\
		\midrule
		Inversiones liquidadas & Inversiones 1 a 15 & Capital devuelto; las
		inversiones 5, 10 y 15 se liquidan 20 días tarde y pagan interés
		moratorio. \\
		Inversiones vencidas sin liquidar & Inversiones 16 a 25 & Los dos
		últimos periodos de cada una quedan sin pago (20 periodos), lo que
		genera adeudos de interés y de capital. \\
		Inversiones vigentes & Inversiones 26 a 40 & Periodos ya pagados y 144
		periodos que aún no vencen. \\
		Inversiones sin aval & 13 inversiones & Llave foránea
		\nem{aval\us num} nula. \\
		Avales sin inversiones & 5 avales & Mínimo cero de la relación
		\nem{aval} -- \nem{inversion}. \\
		Clientes con varias inversiones & 15 clientes & Consultas agrupadas
		por cliente. \\
		Pagos tardíos & 37 pagos de 10 inversiones & Pagos realizados 10 días
		después de la fecha de pago, con interés moratorio. \\
		Pagos parciales & 3 periodos (inversiones 14, 21 y 35) & Un periodo
		cubierto con dos pagos: 60\,\% en la fecha de pago y el resto 15 días
		después, con interés moratorio. \\
		Plazos y montos & 40 inversiones & Plazos de 6, 12, 15, 18 y 24
		meses; capitales de \$50{,}000 a \$500{,}000; tasas de 0.85\,\% a
		1.25\,\% mensual y moratorias de 2.5\,\% y 3.0\,\%. \\
		Distribución en el tiempo & 2024 a 2026 & Suscripciones: 18 en 2024,
		17 en 2025 y 5 en 2026. Pagos: 74 en 2024, 237 en 2025 y 167 en
		2026. \\
		\bottomrule
	\end{tabular}
\end{table}

\begin{table}[H]
	\centering
	\small
	\caption{Volumen de los datos de prueba}
	\label{tab:volumen-prueba}
	\begin{tabular}{@{}l r L{9.5cm}@{}}
		\toprule
		\textbf{Tabla} & \textbf{Registros} & \textbf{Origen} \\
		\midrule
		\nem{empresa} & 1 & Inserción directa \\
		\nem{sucursal} & 15 & Inserción directa \\
		\nem{cliente} & 25 & Inserción directa \\
		\nem{aval} & 15 & Inserción directa \\
		\nem{inversion} & 40 & Inserción directa \\
		\nem{interes} & 639 & Procedimiento \nem{generar\us tabla\us interes},
		un periodo por cada mes de plazo \\
		\nem{liquidacion} & 15 & Derivado de \nem{inversion} \\
		\nem{pago} & 478 & Derivado de \nem{interes}: 435 completos, 37
		tardíos y 6 de periodos parciales \\
		\nem{bitacora} & 0 & La carga inicial no se registra \\
		\bottomrule
	\end{tabular}
\end{table}

\paragraph{Forma de la carga.} El DBA realiza la carga sobre las tablas vacías y antes de crear los triggers de bitácora del capítulo~\ref{cap:bitacora}. La carga inicial no es una operación de los usuarios del sistema, por lo que no se audita. Como las tablas están vacías, las llaves sustitutas se generan del 1 al $N$ en el orden de inserción, y las inversiones referencian a sus sucursales, clientes y avales por esos números. Las inserciones de las tablas independientes y de las inversiones son idénticas en ambos gestores. La moneda toma el valor por defecto (\nem{MXN}).

Los periodos de interés no se insertan como valores literales: se generan con el procedimiento \nem{generar\us tabla\us interes}, el mismo que usa el ejecutivo en la operación, de modo que la carga sigue las reglas del sistema. Los pagos y las liquidaciones tampoco se capturan uno por uno. Se derivan de los periodos y de las inversiones mediante \nem{INSERT \ldots{} SELECT} con reglas fijas, lo que mantiene los montos consistentes y la carga reproducible. Esta es la única parte que difiere entre gestores, por la aritmética de fechas. Las liquidaciones se insertan antes que los pagos, porque la regla de los periodos sin pago identifica a las inversiones vencidas como aquellas que no tienen liquidación.

\subsection{Datos comunes}

\begin{verbatim}
	-- Ejecutado por dlandero en la base de datos inversiones
	INSERT INTO empresa (empresa_rfc, empresa_nombre, empresa_direccion)
	VALUES ('IGO150312AB1', 'Inversiones del Golfo, S.A. de C.V.',
	        'Av. Independencia 1250, Centro, Veracruz, Ver.');

	INSERT INTO sucursal (empresa_rfc, sucursal_nombre, sucursal_ciudad,
	                      sucursal_direccion)
	VALUES
	  ('IGO150312AB1', 'Sucursal Centro', 'Veracruz',
	   'Av. Independencia 1250, Veracruz'),
	  ('IGO150312AB1', 'Sucursal Costa Verde', 'Boca del Río',
	   'Blvd. Manuel Ávila Camacho 3010, Boca del Río'),
	  ('IGO150312AB1', 'Sucursal Xalapa Centro', 'Xalapa',
	   'Av. Xalapa 455, Xalapa'),
	  ('IGO150312AB1', 'Sucursal Córdoba', 'Córdoba',
	   'Calle 3 número 210, Córdoba'),
	  ('IGO150312AB1', 'Sucursal Orizaba', 'Orizaba',
	   'Av. Oriente 6 número 880, Orizaba'),
	  ('IGO150312AB1', 'Sucursal Coatzacoalcos', 'Coatzacoalcos',
	   'Av. Zaragoza 512, Coatzacoalcos'),
	  ('IGO150312AB1', 'Sucursal Poza Rica', 'Poza Rica',
	   'Blvd. Adolfo Ruiz Cortines 1401, Poza Rica'),
	  ('IGO150312AB1', 'Sucursal Tuxpan', 'Tuxpan',
	   'Av. Juárez 37, Tuxpan'),
	  ('IGO150312AB1', 'Sucursal Puebla Angelópolis', 'Puebla',
	   'Blvd. Atlixcáyotl 5320, Puebla'),
	  ('IGO150312AB1', 'Sucursal Ciudad de México Reforma', 'Ciudad de México',
	   'Paseo de la Reforma 222, Ciudad de México'),
	  ('IGO150312AB1', 'Sucursal Monterrey San Pedro', 'Monterrey',
	   'Av. Vasconcelos 150, Monterrey'),
	  ('IGO150312AB1', 'Sucursal Guadalajara Providencia', 'Guadalajara',
	   'Av. Providencia 2577, Guadalajara'),
	  ('IGO150312AB1', 'Sucursal Mérida Norte', 'Mérida',
	   'Calle 60 número 299, Mérida'),
	  ('IGO150312AB1', 'Sucursal Oaxaca Centro', 'Oaxaca',
	   'Av. Juárez 703, Oaxaca'),
	  ('IGO150312AB1', 'Sucursal Villahermosa Tabasco 2000', 'Villahermosa',
	   'Prol. Paseo Usumacinta 1001, Villahermosa');

	INSERT INTO cliente (cliente_rfc, cliente_nombre, cliente_apellido,
	                     cliente_direccion, cliente_telefono)
	VALUES
	  ('HEGJ720604HM3', 'José Luis', 'Hernández García',
	   'Calle Hidalgo 100, Veracruz', '2291000000'),
	  ('LOMM791107RZ6', 'María Fernanda', 'López Martínez',
	   'Calle Morelos 113, Boca del Río', '2291037291'),
	  ('RASC860410YL9', 'Carlos', 'Ramírez Sánchez',
	   'Calle Allende 126, Xalapa', '2291074582'),
	  ('TOGA680913FY2', 'Ana Sofía', 'Torres Gómez',
	   'Calle Juárez 139, Córdoba', '2291111873'),
	  ('FODJ750216NK5', 'Jorge', 'Flores Díaz',
	   'Calle Madero 152, Orizaba', '2291149164'),
	  ('MOCL820719WX8', 'Lucía', 'Morales Cruz',
	   'Calle Hidalgo 165, Coatzacoalcos', '2291186455'),
	  ('REOM891222DJ1', 'Miguel Ángel', 'Reyes Ortiz',
	   'Calle Morelos 178, Poza Rica', '2291223746'),
	  ('JIRD710525LW4', 'Daniela', 'Jiménez Ruiz',
	   'Calle Allende 191, Tuxpan', '2291261037'),
	  ('VACR781001UH7', 'Roberto', 'Vázquez Castillo',
	   'Calle Juárez 204, Puebla', '2291298328'),
	  ('MEAP850304BV0', 'Patricia', 'Mendoza Aguilar',
	   'Calle Madero 217, Ciudad de México', '2291335619'),
	  ('GURF670807JG3', 'Fernando', 'Gutiérrez Romero',
	   'Calle Hidalgo 230, Monterrey', '2291372910'),
	  ('CAHG740110SU6', 'Gabriela', 'Castro Herrera',
	   'Calle Morelos 243, Guadalajara', '2291410201'),
	  ('CAMR810613ZF9', 'Ricardo', 'Chávez Medina',
	   'Calle Allende 256, Mérida', '2291447492'),
	  ('VARA881116GT2', 'Alejandra', 'Vargas Ríos',
	   'Calle Juárez 269, Oaxaca', '2291484783'),
	  ('DOPE700419PE5', 'Eduardo', 'Domínguez Peña',
	   'Calle Madero 282, Villahermosa', '2291522074'),
	  ('SAIV770922XS8', 'Verónica', 'Salazar Ibarra',
	   'Calle Hidalgo 295, Veracruz', '2291559365'),
	  ('NACA840225ED1', 'Arturo', 'Navarro Cortés',
	   'Calle Morelos 308, Boca del Río', '2291596656'),
	  ('ROFC660701MR4', 'Claudia', 'Rojas Fuentes',
	   'Calle Allende 321, Xalapa', '2291633947'),
	  ('GULH731204VC7', 'Héctor', 'Guzmán Luna',
	   'Calle Juárez 334, Córdoba', '2291671238'),
	  ('EIVM800507CP0', 'Mónica', 'Espinoza Vega',
	   'Calle Madero 347, Orizaba', '2291708529'),
	  ('SAPF871010KB3', 'Francisco', 'Sandoval Pacheco',
	   'Calle Hidalgo 360, Coatzacoalcos', '2291745820'),
	  ('CAST690313TN6', 'Teresa', 'Cabrera Soto',
	   'Calle Morelos 373, Poza Rica', '2291783111'),
	  ('COLR760816AA9', 'Raúl', 'Contreras Lara',
	   'Calle Allende 386, Tuxpan', '2291820402'),
	  ('MORS830119HM2', 'Silvia', 'Montes Rosales',
	   'Calle Juárez 399, Puebla', '2291857693'),
	  ('NUVA650622RZ5', 'Andrés', 'Núñez Valdez',
	   'Calle Madero 412, Ciudad de México', '2291894984');

	INSERT INTO aval (aval_rfc, aval_nombre, aval_apellido, aval_direccion)
	VALUES
	  ('PEZM630913ED0', 'Manuel', 'Pérez Zamora',
	   'Av. Revolución 200, Córdoba'),
	  ('DECL700216MR3', 'Laura', 'Delgado Campos',
	   'Av. Insurgentes 217, Orizaba'),
	  ('AOBE770719VC6', 'Ernesto', 'Acosta Bravo',
	   'Av. Constitución 234, Coatzacoalcos'),
	  ('POGR591222CP9', 'Rosa María', 'Ponce Galván',
	   'Av. Reforma 251, Poza Rica'),
	  ('OOTG660525KB2', 'Guillermo', 'Ochoa Trejo',
	   'Av. Universidad 268, Tuxpan'),
	  ('SOMB731001TN5', 'Beatriz', 'Solís Mejía',
	   'Av. Revolución 285, Puebla'),
	  ('CAQA800304AA8', 'Alberto', 'Carrillo Quiroz',
	   'Av. Insurgentes 302, Ciudad de México'),
	  ('BALL620807HM1', 'Leticia', 'Bautista Lozano',
	   'Av. Constitución 319, Monterrey'),
	  ('MOAO690110RZ4', 'Óscar', 'Molina Arias',
	   'Av. Reforma 336, Guadalajara'),
	  ('FITE760613YL7', 'Elena', 'Figueroa Tapia',
	   'Av. Universidad 353, Mérida'),
	  ('VARS581116FY0', 'Sergio', 'Valencia Robles',
	   'Av. Revolución 370, Oaxaca'),
	  ('EUDA650419NK3', 'Adriana', 'Esquivel Duarte',
	   'Av. Insurgentes 387, Villahermosa'),
	  ('CEMJ720922WX6', 'Javier', 'Cervantes Mora',
	   'Av. Constitución 404, Veracruz'),
	  ('BECN790225DJ9', 'Norma', 'Beltrán Cuevas',
	   'Av. Reforma 421, Boca del Río'),
	  ('AESR610701LW2', 'Rubén', 'Avendaño Silva',
	   'Av. Universidad 438, Xalapa');

	INSERT INTO inversion (pagare_num, sucursal_num, cliente_num, aval_num,
	                       inversion_capital, inversion_tasa,
	                       inversion_tasa_moratoria, inversion_dia_pago,
	                       pagare_fecha, pagare_jurisdiccion,
	                       inversion_fecha_vencimiento)
	VALUES
	  ('0001/2024',  1,  1,    1,  50000.00, 0.8500, 3.0000, 10,
	   '2024-01-01', 'Veracruz', '2024-07-10'),
	  ('0002/2024',  2,  2,    2,  75000.00, 0.9500, 2.5000, 19,
	   '2024-02-06', 'Boca del Río', '2025-02-19'),
	  ('0003/2024',  3,  3, NULL, 100000.00, 1.0500, 3.0000, 28,
	   '2024-03-11', 'Xalapa', '2025-09-28'),
	  ('0004/2024',  4,  4,    4, 150000.00, 1.1500, 2.5000,  9,
	   '2024-04-16', 'Córdoba', '2024-10-09'),
	  ('0005/2024',  5,  5,    5, 200000.00, 1.2500, 3.0000, 18,
	   '2024-05-21', 'Orizaba', '2025-05-18'),
	  ('0006/2024',  6,  6, NULL, 250000.00, 0.8500, 2.5000, 27,
	   '2024-06-26', 'Coatzacoalcos', '2025-12-27'),
	  ('0008/2024',  7,  7,    7, 300000.00, 0.9500, 3.0000,  8,
	   '2024-07-04', 'Poza Rica', '2025-01-08'),
	  ('0009/2024',  8,  8,    8, 500000.00, 1.0500, 2.5000, 17,
	   '2024-08-09', 'Tuxpan', '2025-08-17'),
	  ('0011/2024',  9,  9, NULL,  50000.00, 1.1500, 3.0000, 26,
	   '2024-09-14', 'Puebla', '2026-03-26'),
	  ('0013/2024', 10, 10,   10,  75000.00, 1.2500, 2.5000,  7,
	   '2024-10-19', 'Ciudad de México', '2025-04-07'),
	  ('0016/2024', 11, 11,   11, 100000.00, 0.8500, 3.0000, 16,
	   '2024-11-24', 'Monterrey', '2025-11-16'),
	  ('0017/2024', 12, 12, NULL, 150000.00, 0.9500, 2.5000, 25,
	   '2024-12-02', 'Guadalajara', '2026-06-25'),
	  ('0001/2025', 13, 13,   13, 200000.00, 1.0500, 3.0000,  6,
	   '2025-01-07', 'Mérida', '2025-07-06'),
	  ('0003/2025', 14, 14,   14, 250000.00, 1.1500, 2.5000, 15,
	   '2025-02-12', 'Oaxaca', '2026-02-15'),
	  ('0007/2025', 15, 15, NULL, 300000.00, 1.2500, 3.0000, 24,
	   '2025-03-17', 'Villahermosa', '2026-09-24'),
	  ('0007/2024',  1, 16,    1, 500000.00, 0.8500, 2.5000,  5,
	   '2024-07-02', 'Veracruz', '2025-07-05'),
	  ('0010/2024',  2, 17,    2,  50000.00, 0.9500, 3.0000, 14,
	   '2024-08-09', 'Boca del Río', '2025-11-14'),
	  ('0012/2024',  3, 18, NULL,  75000.00, 1.0500, 2.5000, 23,
	   '2024-09-16', 'Xalapa', '2025-09-23'),
	  ('0014/2024',  4, 19,    4, 100000.00, 1.1500, 3.0000,  4,
	   '2024-10-23', 'Córdoba', '2026-01-04'),
	  ('0015/2024',  5, 20,    5, 150000.00, 1.2500, 2.5000, 13,
	   '2024-11-03', 'Orizaba', '2025-11-13'),
	  ('0018/2024',  6, 21, NULL, 200000.00, 0.8500, 3.0000, 22,
	   '2024-12-10', 'Coatzacoalcos', '2026-03-22'),
	  ('0002/2025',  7, 22,    7, 250000.00, 0.9500, 2.5000,  3,
	   '2025-01-17', 'Poza Rica', '2026-01-03'),
	  ('0004/2025',  8, 23,    8, 300000.00, 1.0500, 3.0000, 12,
	   '2025-02-24', 'Tuxpan', '2026-05-12'),
	  ('0006/2025',  9, 24, NULL, 500000.00, 1.1500, 2.5000, 21,
	   '2025-03-04', 'Puebla', '2026-03-21'),
	  ('0009/2025', 10, 25,   10,  50000.00, 1.2500, 3.0000,  2,
	   '2025-04-11', 'Ciudad de México', '2026-07-02'),
	  ('0005/2025', 11,  1,   11,  75000.00, 0.8500, 2.5000, 11,
	   '2025-03-01', 'Monterrey', '2027-03-11'),
	  ('0008/2025', 12,  2, NULL, 100000.00, 0.9500, 3.0000, 20,
	   '2025-04-05', 'Guadalajara', '2027-04-20'),
	  ('0010/2025', 13,  3,   13, 150000.00, 1.0500, 2.5000,  1,
	   '2025-05-09', 'Mérida', '2027-05-01'),
	  ('0011/2025', 14,  4,   14, 200000.00, 1.1500, 3.0000, 10,
	   '2025-06-13', 'Oaxaca', '2027-06-10'),
	  ('0012/2025', 15,  5, NULL, 250000.00, 1.2500, 2.5000, 19,
	   '2025-07-17', 'Villahermosa', '2027-01-19'),
	  ('0013/2025',  1,  6,    1, 300000.00, 0.8500, 3.0000, 28,
	   '2025-08-21', 'Veracruz', '2027-08-28'),
	  ('0014/2025',  2,  7,    2, 500000.00, 0.9500, 2.5000,  9,
	   '2025-09-25', 'Boca del Río', '2027-03-09'),
	  ('0015/2025',  3,  8, NULL,  50000.00, 1.0500, 3.0000, 18,
	   '2025-10-02', 'Xalapa', '2027-10-18'),
	  ('0016/2025',  4,  9,    4,  75000.00, 1.1500, 2.5000, 27,
	   '2025-11-06', 'Córdoba', '2027-05-27'),
	  ('0017/2025',  5, 10,    5, 100000.00, 1.2500, 3.0000,  8,
	   '2025-12-10', 'Orizaba', '2027-12-08'),
	  ('0001/2026',  6, 11, NULL, 150000.00, 0.8500, 2.5000, 17,
	   '2026-01-14', 'Coatzacoalcos', '2027-07-17'),
	  ('0002/2026',  7, 12,    7, 200000.00, 0.9500, 3.0000, 26,
	   '2026-02-18', 'Poza Rica', '2028-02-26'),
	  ('0003/2026',  8, 13,    8, 250000.00, 1.0500, 2.5000,  7,
	   '2026-03-22', 'Tuxpan', '2027-09-07'),
	  ('0004/2026',  9, 14, NULL, 300000.00, 1.1500, 3.0000, 16,
	   '2026-04-26', 'Puebla', '2028-04-16'),
	  ('0005/2026', 10, 15,   10, 500000.00, 1.2500, 2.5000, 25,
	   '2026-05-03', 'Ciudad de México', '2027-11-25');

	-- Tabla de interés de cada inversión
	CALL generar_tabla_interes(1);   CALL generar_tabla_interes(2);
	CALL generar_tabla_interes(3);   CALL generar_tabla_interes(4);
	CALL generar_tabla_interes(5);   CALL generar_tabla_interes(6);
	CALL generar_tabla_interes(7);   CALL generar_tabla_interes(8);
	CALL generar_tabla_interes(9);   CALL generar_tabla_interes(10);
	CALL generar_tabla_interes(11);  CALL generar_tabla_interes(12);
	CALL generar_tabla_interes(13);  CALL generar_tabla_interes(14);
	CALL generar_tabla_interes(15);  CALL generar_tabla_interes(16);
	CALL generar_tabla_interes(17);  CALL generar_tabla_interes(18);
	CALL generar_tabla_interes(19);  CALL generar_tabla_interes(20);
	CALL generar_tabla_interes(21);  CALL generar_tabla_interes(22);
	CALL generar_tabla_interes(23);  CALL generar_tabla_interes(24);
	CALL generar_tabla_interes(25);  CALL generar_tabla_interes(26);
	CALL generar_tabla_interes(27);  CALL generar_tabla_interes(28);
	CALL generar_tabla_interes(29);  CALL generar_tabla_interes(30);
	CALL generar_tabla_interes(31);  CALL generar_tabla_interes(32);
	CALL generar_tabla_interes(33);  CALL generar_tabla_interes(34);
	CALL generar_tabla_interes(35);  CALL generar_tabla_interes(36);
	CALL generar_tabla_interes(37);  CALL generar_tabla_interes(38);
	CALL generar_tabla_interes(39);  CALL generar_tabla_interes(40);
\end{verbatim}

\subsection{Liquidaciones y pagos}

Las reglas para derivar las liquidaciones y los pagos son las mismas en ambos gestores:

\begin{itemize}
	\item Las inversiones 1 a 15 se liquidan en su fecha de vencimiento. Las que tienen un número múltiplo de 5 se liquidan 20 días después y pagan interés moratorio sobre el capital por esos 20 días.
	\item Solo se pagan los periodos cuya fecha final es anterior o igual a la fecha de carga. En las inversiones vencidas sin liquidar, los periodos	que terminan en sus dos últimos meses quedan sin pago.
	\item En las inversiones con número múltiplo de 4, los periodos que	terminan en marzo, junio, septiembre o diciembre se pagan 10 días tarde,	con interés moratorio sobre el interés del periodo.
	\item En las inversiones con número múltiplo de 7, el periodo que termina en junio, si no es tardío, se paga en dos partes: 60\,\% en la fecha de pago y el resto 15 días después, con interés moratorio sobre el resto.
	\item Los demás periodos se pagan completos en su fecha final.
\end{itemize}

El interés moratorio se calcula de forma proporcional a los días de atraso, sobre la tasa moratoria mensual de la inversión y meses de 30 días.

\subsubsection{PostgreSQL}

\begin{verbatim}
	INSERT INTO liquidacion (inversion_num, liquidacion_fecha,
	                         liquidacion_capital, liquidacion_mora)
	SELECT inversion_num,
	       CASE WHEN inversion_num % 5 = 0
	            THEN inversion_fecha_vencimiento + 20
	            ELSE inversion_fecha_vencimiento
	       END,
	       inversion_capital,
	       CASE WHEN inversion_num % 5 = 0
	            THEN ROUND(inversion_capital * inversion_tasa_moratoria
	                       / 100 * 20 / 30, 2)
	            ELSE 0
	       END
	FROM inversion
	WHERE inversion_num BETWEEN 1 AND 15;

	INSERT INTO pago (inversion_num, interes_fecha_inicial, pago_fecha,
	                  pago_interes, pago_mora)
	WITH periodo AS (
	  SELECT i.inversion_num, i.interes_fecha_inicial,
	         i.interes_fecha_final AS fin, i.interes_monto AS monto,
	         v.inversion_tasa_moratoria AS tasa_mor,
	         (v.inversion_num % 4 = 0 AND
	          EXTRACT(MONTH FROM i.interes_fecha_final) IN (3, 6, 9, 12))
	           AS tardio,
	         (v.inversion_num % 7 = 0 AND
	          EXTRACT(MONTH FROM i.interes_fecha_final) = 6) AS parcial
	  FROM interes i
	    JOIN inversion v ON v.inversion_num = i.inversion_num
	  WHERE i.interes_fecha_final <= DATE '2026-10-28'
	    AND NOT (v.inversion_fecha_vencimiento < DATE '2026-10-28'
	             AND NOT EXISTS (SELECT 1 FROM liquidacion l
	                             WHERE l.inversion_num = v.inversion_num)
	             AND i.interes_fecha_final >
	                 v.inversion_fecha_vencimiento - INTERVAL '2 months')
	)
	-- Pago completo en la fecha de pago
	SELECT inversion_num, interes_fecha_inicial, fin, monto, 0
	FROM periodo
	WHERE NOT tardio AND NOT parcial
	UNION ALL
	-- Pago tardío: 10 días de atraso
	SELECT inversion_num, interes_fecha_inicial, fin + 10, monto,
	       ROUND(monto * tasa_mor / 100 * 10 / 30, 2)
	FROM periodo
	WHERE tardio
	UNION ALL
	-- Pago parcial: 60 % en la fecha de pago
	SELECT inversion_num, interes_fecha_inicial, fin,
	       ROUND(monto * 0.6, 2), 0
	FROM periodo
	WHERE parcial AND NOT tardio
	UNION ALL
	-- Resto del pago parcial: 15 días de atraso
	SELECT inversion_num, interes_fecha_inicial, fin + 15,
	       monto - ROUND(monto * 0.6, 2),
	       ROUND((monto - ROUND(monto * 0.6, 2)) * tasa_mor / 100 * 15 / 30, 2)
	FROM periodo
	WHERE parcial AND NOT tardio;
\end{verbatim}

\subsubsection{MySQL}

\begin{verbatim}
	INSERT INTO liquidacion (inversion_num, liquidacion_fecha,
	                         liquidacion_capital, liquidacion_mora)
	SELECT inversion_num,
	       CASE WHEN MOD(inversion_num, 5) = 0
	            THEN DATE_ADD(inversion_fecha_vencimiento, INTERVAL 20 DAY)
	            ELSE inversion_fecha_vencimiento
	       END,
	       inversion_capital,
	       CASE WHEN MOD(inversion_num, 5) = 0
	            THEN ROUND(inversion_capital * inversion_tasa_moratoria
	                       / 100 * 20 / 30, 2)
	            ELSE 0
	       END
	FROM inversion
	WHERE inversion_num BETWEEN 1 AND 15;

	INSERT INTO pago (inversion_num, interes_fecha_inicial, pago_fecha,
	                  pago_interes, pago_mora)
	WITH periodo AS (
	  SELECT i.inversion_num, i.interes_fecha_inicial,
	         i.interes_fecha_final AS fin, i.interes_monto AS monto,
	         v.inversion_tasa_moratoria AS tasa_mor,
	         (MOD(v.inversion_num, 4) = 0 AND
	          MONTH(i.interes_fecha_final) IN (3, 6, 9, 12)) AS tardio,
	         (MOD(v.inversion_num, 7) = 0 AND
	          MONTH(i.interes_fecha_final) = 6) AS parcial
	  FROM interes i
	    JOIN inversion v ON v.inversion_num = i.inversion_num
	  WHERE i.interes_fecha_final <= DATE '2026-10-28'
	    AND NOT (v.inversion_fecha_vencimiento < DATE '2026-10-28'
	             AND NOT EXISTS (SELECT 1 FROM liquidacion l
	                             WHERE l.inversion_num = v.inversion_num)
	             AND i.interes_fecha_final >
	                 DATE_SUB(v.inversion_fecha_vencimiento, INTERVAL 2 MONTH))
	)
	-- Pago completo en la fecha de pago
	SELECT inversion_num, interes_fecha_inicial, fin, monto, 0
	FROM periodo
	WHERE NOT tardio AND NOT parcial
	UNION ALL
	-- Pago tardío: 10 días de atraso
	SELECT inversion_num, interes_fecha_inicial,
	       DATE_ADD(fin, INTERVAL 10 DAY), monto,
	       ROUND(monto * tasa_mor / 100 * 10 / 30, 2)
	FROM periodo
	WHERE tardio
	UNION ALL
	-- Pago parcial: 60 % en la fecha de pago
	SELECT inversion_num, interes_fecha_inicial, fin,
	       ROUND(monto * 0.6, 2), 0
	FROM periodo
	WHERE parcial AND NOT tardio
	UNION ALL
	-- Resto del pago parcial: 15 días de atraso
	SELECT inversion_num, interes_fecha_inicial,
	       DATE_ADD(fin, INTERVAL 15 DAY),
	       monto - ROUND(monto * 0.6, 2),
	       ROUND((monto - ROUND(monto * 0.6, 2)) * tasa_mor / 100 * 15 / 30, 2)
	FROM periodo
	WHERE parcial AND NOT tardio;
\end{verbatim}

\subsection{Resultado de la carga}

La consulta siguiente, igual en ambos gestores, cuenta los registros de cada tabla para compararlos con la tabla~\ref{tab:volumen-prueba}.

\begin{verbatim}
	SELECT 'empresa' AS tabla, COUNT(*) AS registros FROM empresa
	UNION ALL SELECT 'sucursal',    COUNT(*) FROM sucursal
	UNION ALL SELECT 'cliente',     COUNT(*) FROM cliente
	UNION ALL SELECT 'aval',        COUNT(*) FROM aval
	UNION ALL SELECT 'inversion',   COUNT(*) FROM inversion
	UNION ALL SELECT 'interes',     COUNT(*) FROM interes
	UNION ALL SELECT 'liquidacion', COUNT(*) FROM liquidacion
	UNION ALL SELECT 'pago',        COUNT(*) FROM pago
	UNION ALL SELECT 'bitacora',    COUNT(*) FROM bitacora;
\end{verbatim}

\captura{pg-datos-conteo.png}{Registros cargados en PostgreSQL}{pg-datos-conteo}
\captura{mysql-datos-conteo.png}{Registros cargados en MySQL}{mysql-datos-conteo}

\section{Pruebas de alcance}
\label{sec:pruebas-alcance}
Las pruebas se ejecutan después de crear los triggers de bitácora del capítulo~\ref{cap:bitacora}, para que las operaciones permitidas queden auditadas.

Las pruebas de alcance comprueban que cada rol puede ejecutar las operaciones que le asigna la matriz de permisos de la sección~\ref{sec:roles-matriz} y que el gestor rechaza las demás. Cada usuario ejecuta al menos una operación permitida y una denegada. Las operaciones son las que realizaría el puesto en la operación diaria, de modo que las pruebas permitidas quedan como registros reales del sistema: un cliente, una inversión con su tabla de interés y un pago. Las sentencias son iguales en ambos gestores; solo cambian la conexión y el texto del error. La tabla~\ref{tab:pruebas-alcance} resume las pruebas.

\begin{table}[H]
	\centering
	\small
	\caption{Pruebas de alcance por rol}
	\label{tab:pruebas-alcance}
	\begin{tabular}{@{}L{3.2cm} L{8.2cm} L{4.6cm}@{}}
		\toprule
		\textbf{Usuario (rol)} & \textbf{Operación} & \textbf{Resultado	esperado} \\
		\midrule
		\nem{dmejia} \newline (\nem{rol\us gerente}) & Actualizar la dirección de una sucursal & Permitida \\
		& Registrar un pago & Denegada: sin privilegio \\
		& Eliminar una inversión con periodos de interés & Denegada: llave foránea (nota 1 de la matriz) \\
		\midrule
		\nem{pzuniga} \newline (\nem{rol\us ejecutivo}) & Registrar un cliente y una inversión, generar su tabla de interés y actualizar el teléfono del cliente & Permitida \\
		& Modificar el RFC del cliente & Denegada: privilegio por columna \\
		& Insertar un periodo de interés directamente & Denegada: sin privilegio \\
		\midrule
		\nem{xzaleta} \newline (\nem{rol\us tesoreria}) & Registrar el pago	tardío de un periodo vencido & Permitida \\
		& Modificar el capital de una inversión & Denegada: sin privilegio \\
		\midrule
		\nem{enovoa} \newline (\nem{rol\us auditor}) & Consultar los pagos con interés moratorio & Permitida \\
		& Eliminar registros de la bitácora & Denegada: sin privilegio \\
		\midrule
		\nem{respaldo} \newline (\nem{rol\us respaldo}) & Generar un respaldo lógico de la base de datos & Permitida \\
		& Eliminar pagos & Denegada: sin privilegio \\
		\bottomrule
	\end{tabular}
\end{table}

\subsection{Sentencias de prueba}

\paragraph{Gerente.} La eliminación de la inversión 40 es una operación que el gerente tiene permitida, pero el gestor la rechaza porque la inversión ya tiene periodos de interés.

\begin{verbatim}
	-- Permitida
	UPDATE sucursal
	SET sucursal_direccion = 'Av. Paseo Tabasco 1203, Villahermosa'
	WHERE sucursal_num = 15;

	-- Denegada: sin privilegio
	INSERT INTO pago (inversion_num, interes_fecha_inicial, pago_fecha,
	                  pago_interes)
	VALUES (16, '2025-05-05', '2026-10-28', 4250.00);

	-- Denegada: llave foránea
	DELETE FROM inversion WHERE inversion_num = 40;
\end{verbatim}

\paragraph{Ejecutivo.} El ejecutivo registra una nueva clienta y su inversión, que recibe el número 41, y genera su tabla de interés de 12 periodos mediante el procedimiento.

\begin{verbatim}
	-- Permitidas
	INSERT INTO cliente (cliente_rfc, cliente_nombre, cliente_apellido,
	                     cliente_direccion, cliente_telefono)
	VALUES ('OESM920314HT5', 'Mariana', 'Ortega Salinas',
	        'Calle Enríquez 58, Xalapa', '2281234567');

	INSERT INTO inversion (pagare_num, sucursal_num, cliente_num, aval_num,
	                       inversion_capital, inversion_tasa,
	                       inversion_tasa_moratoria, inversion_dia_pago,
	                       pagare_fecha, pagare_jurisdiccion,
	                       inversion_fecha_vencimiento)
	VALUES ('0006/2026', 3, 26, NULL, 120000.00, 1.0000, 2.5000, 15,
	        '2026-10-28', 'Xalapa', '2027-10-15');

	CALL generar_tabla_interes(41);

	UPDATE cliente SET cliente_telefono = '2287654321'
	WHERE cliente_num = 26;

	-- Denegada: privilegio por columna
	UPDATE cliente SET cliente_rfc = 'OESM920314HT6'
	WHERE cliente_num = 26;

	-- Denegada: sin privilegio
	INSERT INTO interes (inversion_num, interes_fecha_inicial,
	                     interes_fecha_final, interes_monto)
	VALUES (41, '2027-10-15', '2027-11-15', 1200.00);
\end{verbatim}

\paragraph{Tesorería.} La tesorería paga el periodo del 5 de mayo al 5 de junio de 2025 de la inversión 16, que estaba vencido. El interés moratorio corresponde a 510 días de atraso: $4{,}250.00 \times 2.5\mkern 3mu\% \times 510/30 = 1{,}806.25$.


\begin{verbatim}
	-- Permitida
	INSERT INTO pago (inversion_num, interes_fecha_inicial, pago_fecha,
	                  pago_interes, pago_mora)
	VALUES (16, '2025-05-05', '2026-10-28', 4250.00, 1806.25);

	-- Denegada: sin privilegio
	UPDATE inversion SET inversion_capital = 600000.00
	WHERE inversion_num = 16;
\end{verbatim}

\paragraph{Auditor.}

\begin{verbatim}
	-- Permitida
	SELECT inversion_num, COUNT(*) AS pagos, SUM(pago_mora) AS mora
	FROM pago
	WHERE pago_mora > 0
	GROUP BY inversion_num
	ORDER BY inversion_num;

	-- Denegada: sin privilegio
	DELETE FROM bitacora;
\end{verbatim}

\paragraph{Respaldo.} La operación permitida es el respaldo lógico, que se ejecuta con la herramienta de cada gestor (ver las secciones siguientes). La operación denegada es:

\begin{verbatim}
	-- Denegada: sin privilegio
	DELETE FROM pago WHERE inversion_num = 1;
\end{verbatim}

\subsection{PostgreSQL}

Cada usuario se conecta con su propia cuenta y ejecuta sus sentencias:

\begin{verbatim}
	psql -h localhost -p 5432 -U dmejia -d inversiones
	psql -h localhost -p 5432 -U pzuniga -d inversiones
	psql -h localhost -p 5432 -U xzaleta -d inversiones
	psql -h localhost -p 5432 -U enovoa -d inversiones
	psql -h localhost -p 5432 -U respaldo -d inversiones
	
	pg_dump -h localhost -p 5432 -U respaldo -d inversiones -n inversiones -F c -f inversiones.dump
\end{verbatim}

Las operaciones permitidas responden con \nem{UPDATE 1}, \nem{INSERT 0 1}, \nem{CALL} o el resultado de la consulta. Las denegadas por falta de privilegio responden con \nem{ERROR: permission denied for table}, seguido del nombre de la tabla, incluso en la actualización del RFC, porque el ejecutivo no tiene privilegio sobre esa columna. La eliminación de la inversión responde con \nem{ERROR: update or delete on table "inversion" violates foreign key constraint \ldots{} on table "interes"}.

\captura{pg-alcance-gerente.png}{Pruebas de alcance del gerente en PostgreSQL}{pg-alcance-gerente}
\captura{pg-alcance-ejecutivo.png}{Pruebas de alcance del ejecutivo en PostgreSQL}{pg-alcance-ejecutivo}
\captura{pg-alcance-tesoreria.png}{Pruebas de alcance de tesorería en PostgreSQL}{pg-alcance-tesoreria}
\captura{pg-alcance-auditor.png}{Pruebas de alcance del auditor en PostgreSQL}{pg-alcance-auditor}
\captura{pg-alcance-respaldo.png}{Pruebas de alcance del respaldo en PostgreSQL}{pg-alcance-respaldo}

\subsection{MySQL}

\begin{verbatim}
	mysql -h localhost -P 3306 -u dmejia -p inversiones
	mysql -h localhost -P 3306 -u pzuniga -p inversiones
	mysql -h localhost -P 3306 -u xzaleta -p inversiones
	mysql -h localhost -P 3306 -u enovoa -p inversiones
	mysql -h localhost -P 3306 -u respaldo -p inversiones
	
	mysqldump -h localhost -P 3306 -u respaldo -p --single-transaction --routines --triggers inversiones > inversiones.sql
\end{verbatim}

Las operaciones permitidas responden con \nem{Query OK}. Las denegadas por falta de privilegio sobre la tabla responden con el error 1142: \nem{INSERT command denied to user \ldots{} for table 'pago'}. La actualización del RFC responde con el error 1143, propio de los privilegios por columna: \nem{UPDATE command denied to user \ldots{} for column 'cliente\us rfc' in table 'cliente'}. La eliminación de la inversión responde con el error 1451: \nem{Cannot delete or update a parent row: a foreign key constraint fails}.

\captura{mysql-alcance-gerente.png}{Pruebas de alcance del gerente en MySQL}{mysql-alcance-gerente}
\captura{mysql-alcance-ejecutivo.png}{Pruebas de alcance del ejecutivo en MySQL}{mysql-alcance-ejecutivo}
\captura{mysql-alcance-tesoreria.png}{Pruebas de alcance de tesorería en MySQL}{mysql-alcance-tesoreria}
\captura{mysql-alcance-auditor.png}{Pruebas de alcance del auditor en MySQL}{mysql-alcance-auditor}
\captura{mysql-alcance-respaldo.png}{Pruebas de alcance del respaldo en MySQL}{mysql-alcance-respaldo}